% ============================================================
\chapter{Introduction to Ergodic Theory}
\label{ch:ergodic}
% ============================================================

\section{Introduction}

Ergodic theory\index{ergodic theory} studies the long-term statistical
behaviour of dynamical systems. It addresses the fundamental question: do
time averages coincide with space averages? This chapter presents
measure-preserving transformations, the fundamental ergodic theorems, and
their applications.

\begin{intuition}
Consider a gas in a box. Rather than tracking every molecule, we measure
macroscopic quantities (temperature, pressure) that are averages. Boltzmann's
ergodic hypothesis asserts that the time average of an observable along a
trajectory equals its average over the entire phase space. This is the
foundation of statistical mechanics.
\end{intuition}

\section{Measure-Preserving Transformations}

\begin{definition}[Measure-preserving dynamical system]
\index{measure-preserving transformation}
A \textbf{measure-preserving dynamical system} is a tuple
$(X, \mathcal{B}, \mu, T)$ where $(X, \mathcal{B}, \mu)$ is a probability
space and $T : X \to X$ is a measurable transformation that
\textbf{preserves the measure}: $\mu(T^{-1}(B)) = \mu(B)$ for all
$B \in \mathcal{B}$.
\end{definition}

\begin{example}[Irrational circle rotation]
\index{irrational rotation}
Let $T_\alpha : \R/\Z \to \R/\Z$ be defined by
$T_\alpha(x) = x + \alpha \pmod{1}$ with $\alpha \notin \Q$. Lebesgue
measure on $[0,1)$ is invariant under $T_\alpha$.
\end{example}

\begin{example}[Doubling map]
The map $T(x) = 2x \pmod{1}$ on $[0,1)$ preserves Lebesgue measure. It is
a non-invertible endomorphism with entropy $\ln 2$.
\end{example}

\begin{example}[Baker's transformation]
\index{baker's transformation}
The \textbf{baker's transformation} on $[0,1)^2$ is
\[
    B(x, y) = \begin{cases}
    (2x, y/2) & \text{if } 0 \leq x < 1/2, \\
    (2x - 1, (y+1)/2) & \text{if } 1/2 \leq x < 1.
    \end{cases}
\]
It preserves Lebesgue measure and models kneading: horizontal stretching,
vertical compression, and folding.
\end{example}

\section{Ergodicity}

\begin{definition}[Invariant set]
A set $A \in \mathcal{B}$ is \textbf{$T$-invariant} if $T^{-1}(A) = A$
(modulo a null set).
\end{definition}

\begin{definition}[Ergodic transformation]
\index{ergodicity}
$T$ is \textbf{ergodic} if every $T$-invariant set has measure $0$ or $1$.
Equivalently, every measurable $T$-invariant function
($\varphi \circ T = \varphi$ a.e.) is constant a.e.
\end{definition}

\begin{theorem}[Ergodicity of irrational rotations]
The rotation $T_\alpha$ is ergodic with respect to Lebesgue measure if and
only if $\alpha$ is irrational.
\end{theorem}

\begin{proof}
Let $\varphi \in L^2([0,1))$ be invariant: $\varphi(x + \alpha) = \varphi(x)$
a.e. Expanding in Fourier series
$\varphi(x) = \sum_{n \in \Z} c_n e^{2\pi inx}$, invariance gives
$c_n e^{2\pi in\alpha} = c_n$ for all $n$. If $\alpha \notin \Q$, then
$e^{2\pi in\alpha} \neq 1$ for $n \neq 0$, so $c_n = 0$ for all $n \neq 0$.
Thus $\varphi = c_0$ is constant a.e.
\end{proof}

\section{Birkhoff's Ergodic Theorem}

\begin{theorem}[Birkhoff, 1931]
\label{thm:birkhoff}
\index{Birkhoff's ergodic theorem}
Let $(X, \mathcal{B}, \mu, T)$ be a measure-preserving system and
$\varphi \in L^1(X, \mu)$. Then, for $\mu$-almost every $x$:
\[
    \overline{\varphi}(x) \coloneqq \lim_{N \to \infty} \frac{1}{N}
    \sum_{n=0}^{N-1} \varphi(T^n(x))
\]
exists and $\overline{\varphi}$ is $T$-invariant. Moreover,
$\int_X \overline{\varphi}\,\mathrm{d}\mu = \int_X \varphi\,\mathrm{d}\mu$.
\end{theorem}

\begin{corollary}[Ergodic case]
If $T$ is ergodic, then for every $\varphi \in L^1$ and $\mu$-almost every $x$:
\[
    \lim_{N \to \infty} \frac{1}{N} \sum_{n=0}^{N-1} \varphi(T^n(x))
    = \int_X \varphi\,\mathrm{d}\mu.
\]
\textbf{Time averages equal space averages.}
\end{corollary}

\begin{keyformulas}[title=\textbf{Fundamental Ergodic Theorems}]
\begin{itemize}
    \item \textbf{Von Neumann (1932)}: $L^2$ convergence of
    $\frac{1}{N}\sum_{n=0}^{N-1}\varphi \circ T^n$ to the projection onto
    $T$-invariant functions.
    \item \textbf{Birkhoff (1931)}: almost sure convergence (stronger).
    \item \textbf{Kingman (1968)}: subadditive ergodic theorem for sequences
    $(a_n)$ with $a_{m+n} \leq a_m + a_n \circ T^m$.
\end{itemize}
\end{keyformulas}

\begin{example}[Weyl equidistribution]
Applying Birkhoff to $T_\alpha$ with $\varphi = \ind_{[a,b]}$ yields Weyl's
theorem: for $\alpha \notin \Q$,
\[
    \lim_{N \to \infty}
    \frac{\#\{0 \leq n < N : \{n\alpha\} \in [a,b]\}}{N} = b - a.
\]
The sequence $(n\alpha \bmod 1)_{n \geq 0}$ is equidistributed modulo~1.
\end{example}

\section{Mixing}

\begin{definition}[Weak and strong mixing]
\index{mixing}
$T$ is \textbf{weakly mixing} if for all $A, B \in \mathcal{B}$:
\[
    \lim_{N \to \infty} \frac{1}{N} \sum_{n=0}^{N-1}
    \abs{\mu(T^{-n}(A) \cap B) - \mu(A)\mu(B)} = 0.
\]
$T$ is \textbf{(strongly) mixing} if:
\[
    \lim_{n \to \infty} \mu(T^{-n}(A) \cap B) = \mu(A)\mu(B).
\]
\end{definition}

\begin{proposition}[Implications]
\[
    \text{strong mixing} \implies \text{weak mixing} \implies \text{ergodicity}.
\]
Neither implication reverses in general.
\end{proposition}

\begin{example}
The doubling map $T(x) = 2x \pmod{1}$ and the baker's transformation are
strongly mixing. The irrational rotation is ergodic but \textbf{not} mixing
(not even weakly).
\end{example}

\begin{theorem}[Spectral characterisation of mixing]
$T$ is strongly mixing if and only if for all
$\varphi, \psi \in L^2(X, \mu)$:
\[
    \lim_{n \to \infty} \int_X (\varphi \circ T^n) \cdot \psi\,\mathrm{d}\mu
    = \int_X \varphi\,\mathrm{d}\mu \cdot \int_X \psi\,\mathrm{d}\mu.
\]
\end{theorem}

\section{Ergodic Decomposition}

\begin{theorem}[Ergodic decomposition]
\index{ergodic decomposition}
Let $T : X \to X$ preserve $\mu$ on a standard probability space. Then there
exists an (essentially unique) decomposition
\[
    \mu = \int_Y \mu_y\,\mathrm{d}\nu(y)
\]
where $\nu$ is a probability measure on a measurable space $Y$ and each
$\mu_y$ is a $T$-invariant \textbf{ergodic} probability measure.
\end{theorem}

\begin{remark}
Ergodic decomposition means that every measure-preserving system can be
decomposed into irreducible ergodic components. It is the analogue of
decomposing a Hilbert space into irreducible subspaces.
\end{remark}

\section{Metric Entropy}

\begin{definition}[Kolmogorov-Sinai entropy]
\index{Kolmogorov-Sinai entropy}
Let $\mathcal{P} = \{P_1, \ldots, P_k\}$ be a measurable partition of $X$.
The \textbf{entropy} of $\mathcal{P}$ is
$H(\mathcal{P}) = -\sum_{i=1}^k \mu(P_i)\ln\mu(P_i)$. The
\textbf{metric entropy} of $T$ with respect to $\mathcal{P}$ is
\[
    h_\mu(T, \mathcal{P}) = \lim_{n \to \infty} \frac{1}{n}
    H\!\left(\bigvee_{k=0}^{n-1} T^{-k}\mathcal{P}\right),
\]
and the \textbf{Kolmogorov-Sinai entropy} is
$h_\mu(T) = \sup_{\mathcal{P}} h_\mu(T, \mathcal{P})$.
\end{definition}

\begin{theorem}[Kolmogorov-Sinai]
If $\mathcal{P}$ is a generating partition (i.e.,
$\bigvee_{n=0}^{\infty} T^{-n}\mathcal{P}$ generates $\mathcal{B}$ modulo
null sets), then $h_\mu(T) = h_\mu(T, \mathcal{P})$.
\end{theorem}

\section{Applications to Number Theory}

\begin{example}[Continued fractions and the Gauss map]
\index{Gauss map}
The \textbf{Gauss map} $G : (0,1] \to [0,1)$ defined by $G(x) = \{1/x\}$
(fractional part of $1/x$) encodes the continued fraction algorithm. The
Gauss measure $\mathrm{d}\mu_G = \frac{1}{\ln 2}\frac{\mathrm{d}x}{1+x}$
is $G$-invariant and ergodic.
\end{example}

\begin{corollary}[Khintchine's theorem]
For almost every $x \in (0,1)$ with continued fraction expansion
$x = [a_1, a_2, \ldots]$:
\[
    \lim_{n \to \infty} (a_1 a_2 \cdots a_n)^{1/n}
    = \prod_{k=1}^{\infty}
    \left(1 + \frac{1}{k(k+2)}\right)^{\ln k/\ln 2}
    \approx 2.685\ldots
\]
(Khintchine's constant).
\end{corollary}

\begin{example}[L\'evy's theorem]
By the ergodic theorem applied to $G$:
\[
    \lim_{n \to \infty} \frac{1}{n}\ln q_n
    = \frac{\pi^2}{12\ln 2} \approx 1.1866\ldots
\]
where $q_n$ is the denominator of the $n$-th convergent. This gives the
typical rate of approximation by continued fractions.
\end{example}

\section{Exercises}

\begin{exercise}[Ergodicity]
Show that the baker's transformation is ergodic using Fourier coefficients.
\end{exercise}

\begin{exercise}[Poincar\'e recurrence theorem]
Prove that if $T$ preserves $\mu$ and $\mu(A) > 0$, then for $\mu$-almost
every $x \in A$, there exists $n \geq 1$ with $T^n(x) \in A$.
\end{exercise}

\begin{exercise}[Mixing of the doubling map]
Show that $T(x) = 2x \pmod{1}$ is strongly mixing using the spectral
characterisation with the functions $e^{2\pi ikx}$.
\end{exercise}

\begin{exercise}[Entropy]
Compute the metric entropy of the Gauss map $G(x) = \{1/x\}$ with respect
to the Gauss measure. (Answer: $h_{\mu_G}(G) = \pi^2/(6\ln 2)$.)
\end{exercise}

\begin{exercise}[Equidistribution]
Use Birkhoff's theorem to show that the sequence $(\sin(n) \bmod 1)_{n \geq 1}$
is not necessarily equidistributed, and discuss conditions on $\alpha$ for
$(n^\alpha \bmod 1)$ to be equidistributed.
\end{exercise}

\begin{exercise}[Gauss map verification]
Numerically verify Khintchine's constant by computing the geometric mean of
the continued fraction coefficients of random numbers.
\end{exercise}
